\documentclass{zhvt-classic}

% 使用本地源流明体（GenRyuMin2 TW）按文件路径加载，不依赖系统字体名
% 版框（内框／界栏／书口侧栏／鱼尾）改用与版心印章同一印泥棕色：
%   印章图片 yshs.png 实色 = RGB(135,68,52) = HTML 874434
\usepackage{xcolor}
\usepackage{tikz}
% sealbrown 双注册：xcolor 供 TikZ/pgf，L3 供 cls 版框
\definecolor{sealbrown}{HTML}{874434}
\ExplSyntaxOn
\color_set:nnn { sealbrown } { HTML } { 874434 }
\ExplSyntaxOff

\zhvtset{
  grid color       = sealbrown,
  main font        = GenRyuMin2TW-SB,
  main font path   = C:/Users/Administrator/texmf/fonts/opentype/public/GenRyuMin2TW/,
jiazhu font      = GenRyuMin2TW-M,
  jiazhu font path = C:/Users/Administrator/texmf/fonts/opentype/public/GenRyuMin2TW/,
  % 页码框容纳 4 字：正常页「九十七」+前置空 1 格=4 行恰好；
  % 100+ 页「一百零一」共 5 行超出，末字「一/二/三」被框裁掉（2026-10-04 实测），故放宽到 5
  shukou page number chars = 5,
  % 书口书名框高从默认 6 字减到 4 字，把省出的 2 字高拨给 chapter 框，
  % 使 6 字书名+卷名放进 9 字高的 chapter 框后不再压线（2026-10-05）
  shukou title chars = 4,
}

\title{史记}[张丞相列传]
\author{汉·司马迁}
\maker{闲叟漫记·重製}


% ===== 索隐述赞样式（2026-11-08 四次定稿 · 愚大人亲手改定参数）=====
% ⚠ 竖排前提：自定义字体家族必须带 RawFeature={vertical:+vert}；
%   simhei.ttf 加 +vhal 会整体加载失败静默回退，只写 +vert；simfang 同理。
\newCJKfontfamily{\shuzanfont}[Path=C:/Windows/Fonts/,RawFeature={vertical:+vert}]{simhei.ttf}
\newCJKfontfamily{\shuzanbody}[Path=C:/Windows/Fonts/,RawFeature={vertical:+vert}]{simfang.ttf}
% 题头：152x30pt 圆角 6pt 棕底白字；baseline (0,2) 使整框落入界行格子不压版框线；
% 白字节点 x 不补偿、y +8.2pt（像素级墨迹实测 dx/dy≈0，L/R 留白约 20pt）。
\newcommand{\shuzanhead}[1]{%
  \begin{tikzpicture}[baseline={(0,2)}]
    \fill[sealbrown, rounded corners=6pt]
      (-76pt,-15pt) rectangle (76pt,15pt);
    \node[text=white, inner sep=0pt, anchor=center] at (0pt,8.2pt)
      {\shuzanfont\fontsize{23}{23}\selectfont #1};
  \end{tikzpicture}%
}
% ★ 修正书口卷名位置：chapter 框起始位置改回 cls 默认值，避免全局影响。
%   只在每个卷首页通过 \juan 命令里的空命令微调该页位置（安全，不会跨页污染）。
\ExplSyntaxOn
\AtBeginDocument{
  \dim_set:Nn \l__zhvt_shukou_chapter_title_start_dim
    { \l__zhvt_shukou_title_height_dim + \l__zhvt_font_size_dim * 2 }
  \dim_set:Nn \l__zhvt_shukou_page_number_start_dim
    { \l__zhvt_shukou_chapter_title_start_dim + \l__zhvt_shukou_chapter_title_height_dim }
}
\ExplSyntaxOff

\begin{document}

% ★★ 版心（书口）排布覆写：把书名放到鱼尾之下（2026-10-04）
%
%   cls 默认版心自上而下：书名（title 静态框，在鱼尾之上）→ 鱼尾 → 篇名（chapter 动态框）→ 页码
%   本模板改为：          鱼尾（位置不动）→ 书名 → 篇名 → 页码分割线 → 页码
%
%   手法：把 cls 放在 title 静态框里的书名，改放进 chapter 动态框（鱼尾之下）；
%         title 静态框留空。书名取 \l__zhvt_title_str，即 \title{...} 的第一参数，
%         所以换书时只改 \title，不必在此重复写书名。
%
%   ⚠ 三条必须遵守的写法（2026-10-04 实测）：
%     ① 汉字必须用 \hbox{...} 包住。flowfram 会把框内容装进一个宽 14.4mm（= 一个 line_sep
%        = 书口栏宽）的纵向 minipage 并按该宽度折行，而整页靠 PDF /Rotate 90 显示、
%        框内文字本身不变向 —— 不包 \hbox 时「围炉夜话」4 字会被逐字折成 4 行，
%        在版面上就成了横排 4 字并横向溢出版框。
%     ② 竖排下**一个 \hbox 就是一整列**。书名与篇名要同处一列，必须放进同一个 \hbox；
%        分成两个 \hbox 时，后一个会被当成新的一列、甩到隔壁栏去
%        （实测「序」跑到书口栏外侧 14.4mm 的正文列上）。中间的 \quad 在框内是横向空白，
%        显示出来正好是书名与篇名之间的纵向间隔。
%     ③ ★ 内容盒要写成 **与框等宽的定宽盒** `\hbox to \l__zhvt_shukou_chapter_title_height_dim{…\hfil}`。
%        实测：chapter 框是把内容盒沿纵向**居中**摆放的，内容宽度一变，书名就跟着漂：
%        自然宽度的 \hbox（书名 38.4mm）落在 x=157.2mm，加篇名后（57.6mm）落在 x=152.4mm，
%        目录页（无篇名）与正文页的书名就差半格。做成等宽盒后，无论有没有篇名，
%        内容固定从框顶起排，实测两页均落在 x=142.8mm，不再漂。
\ExplSyntaxOn
% 把书口 chapter 框（书名+篇名）起始位置下移 1 字高，避免 4 字卷名顶到鱼尾
% ⚠ cls 的 \__zhvt_calculate_dimensions: 在 \AtBeginDocument 里执行，会重算该 dim；
%   必须把覆写挂到 begin document 之后，否则会被二次覆盖（2026-10-05 实测）。
% (旧方案已移到导言区，直接修改 chapter 框起始位置)
% 卷标题命令 #1=书口篇名(鱼尾) #2=正文卷名（含卷序，卷名内空格用 ~）
%
% ★★ 目录页两项修订（2026-10-05，愚大人指出）★★
%   ① 页码起点统一对齐：
%      原状 —— 目录项写成 \par\hspace{4em}卷名\hspace{4em}页码，页码起点 = 卷名实宽 + 4em。
%      卷名实宽随字数变：「卷一~梁惠王上」6 字 +~、「卷七~离娄上」5 字 +~，
%      差 1 个字位 → 卷七~卷十 4 条的页码整体比其余 10 条高 27.21pt（一个字），看着上下不齐。
%      改法 —— 卷名补全角空格（宽 = 1 字 = 27.215pt）补足到 **7 字符位**再写进目录，
%      14 条页码起点遂全部落在同一条竖线上（起点对齐；末字随页码字数自然参差）。
%      ⚠ 只对目录里的副本 \l_tmpc_tl 补，正文卷名仍用 \l_tmpa_tl 原样 —— 正文像素不受影响。
%   ② 卷名整体上移半字位：前导 \hspace 由 4em 改 3.5em。
%      实测原状目录块在版框内偏下（天头留白 139.6pt / 地脚留白 110.7pt），
%      上移 0.5em = 13.6pt 后上下基本均分，视觉居中。
\NewDocumentCommand { \juan } { mm } {
	\clearpage
	\tl_set:Nn \l__zhvt_chapter_title_tl { #1 }
	\tl_set:Nn \l_tmpa_tl { #2 }
	\tl_set_eq:NN \l_tmpc_tl \l_tmpa_tl
	\int_step_inline:nn { 7 - \str_count:n { #2 } }
	  { \tl_put_right:Nn \l_tmpc_tl { 　 } }
	\tl_set:Nx \l_tmpb_tl { \zhnum { page } }
	\tl_put_left:Nn \l_tmpb_tl { \hspace { 4em } }
	\tl_put_left:Nn \l_tmpb_tl { \l_tmpc_tl }
	\tl_put_left:Nn \l_tmpb_tl { \par \hspace { 3.5em } }
	\addtocontents { toc } { \l_tmpb_tl }
	\skip_horizontal:n { \l__zhvt_font_size_dim * 4 }
	\l_tmpa_tl
	\par
	\thispagestyle { empty }
}
\cs_set_protected:Npn \__zhvt_set_shukou_content: {
	% ★ 鱼尾上方留白（title 静态框）嵌入闲叟钤印：
	%   整页 /Rotate 90，故此处把图用 \rotatebox{90} 转正为竖章；
	%   尺寸卡在版心条幅内（书内约 7mm 宽 × 33.8mm 高，留余量不压鱼尾）。
	\setstaticcontents * { title } {
		\vbox_to_ht:nn { \l__zhvt_baseline_skip_dim } {
			\vskip 14.2pt
			\hbox_to_wd:nn { \l__zhvt_shukou_title_height_dim } {
				\hfil
				\rotatebox [ origin = c ] { 90 } { \includegraphics [ width = 7mm , height = 33.8mm ] { yshs.png } }
				\hfil
			}
		}
	}
% ★ 版心（书口）统一改法（2026-11-16 愚大人定）：18pt＋框内居中＋书口文字同版框色
%   ① 原 27.2pt 时 7 字长篇名（如「廉颇蔺相如列传」）末字溢出 chapter 框、压住页码；
%      改 \fontsize{18pt}{18pt}，前后 \hfil 令字串在框内居中（不随篇名字数漂移）。
%   ② \vbox{\vskip 8.2pt} 定量下移，抵消缩字号后栏宽方向 8.2pt 上偏（像素实测 y 中心 543.3 = 原版）。
%      ⚠ \raisebox / \vspace* 在 flowfram 框内无效；换字号须同步换补偿：16pt→10.1、19pt→7.4。
%   ③ chapter 与 page~number 均置 \color{sealbrown}（HTML 874434），与版框／界行同色；正文仍黑。
%   ⚠ 页码字位、钤印、鱼尾、界行、正文一律未动；五部合订本同此口径。
	\setdynamiccontents * { chapter }
	{ \color{sealbrown} \vbox { \vskip 8.2pt \hbox to \l__zhvt_shukou_chapter_title_height_dim
		{ \hfil \fontsize{18pt}{18pt}\selectfont \l__zhvt_title_str \quad \l__zhvt_chapter_title_tl \hfil } } }
	\setdynamiccontents * { page~number } { \color{sealbrown} \quad \zhnum { page } }
}
\__zhvt_set_shukou_content:
\ExplSyntaxOff

\maketitle

% 单篇，无多卷目录，不渲染目录页

\juan{张丞相列传}{卷九十六~张丞相列传第三十六}
\hspace*{2em}张丞相苍者。阳武人也。好书律历。秦时为御史。主柱下方书。有罪。亡归。及沛公略地过阳武。苍以客从攻南阳。苍坐法当斩。解衣伏质。身长大。肥白如瓠。时王陵见而怪其美士。乃言沛公。赦勿斩。遂从西入武关。至咸阳。沛公立为汉王。入汉中。还定三秦。陈馀击走常山王张耳。耳归汉。汉乃以张苍为常山守。从淮阴侯击赵。苍得陈馀。赵地已平。汉王以苍为代相。备边寇。已而徙为赵相。相赵王耳。耳卒。相赵王敖。复徙相代王。燕王臧荼反。高祖往击之。苍以代相从攻臧荼有功。以六年中封为北平侯。食邑千二百户。\\
\hspace*{2em}迁为计相。一月。更以列侯为主计四岁。是时萧何为相国。而张苍乃自秦时为柱下史。明习天下图书计籍。苍又善用算律历。故令苍以列侯居相府。领主郡国上计者。黥布反亡。汉立皇子长为淮南王。而张苍相之。十四年。迁为御史大夫。\\
\hspace*{2em}周昌者。沛人也。其从兄曰周苛。秦时皆为泗水卒史。及高祖起沛。击破泗水守监。於是周昌。周苛自卒史从沛公。沛公以周昌为职志。周苛为客。从入关。破秦。沛公立为汉王。以周苛为御史大夫。周昌为中尉。\\
\hspace*{2em}汉王四年。楚围汉王荥阳急。汉王遁出去。而使周苛守荥阳城。楚破荥阳城。欲令周苛将。苛骂曰。﹁若趣降汉王。不然。今为虏矣。﹂项羽怒。亨周苛。於是乃拜周昌为御史大夫。常从击破项籍。以六年中与萧。曹等俱封。封周昌为汾阴侯。周苛子周成以父死事。封为高景侯。\\
\hspace*{2em}昌为人彊力。敢直言。自萧。曹等皆卑下之。昌尝燕时入奏事。高帝方拥戚姬。昌还走。高帝逐得。骑周昌项。问曰。﹁我何如主也。﹂昌仰曰。﹁陛下即桀纣之主也。﹂於是上笑之。然尤惮周昌。及帝欲废太子。而立戚姬子如意为太子。大臣固争之。莫能得。上以留侯策即止。而周昌廷争之彊。上问其说。昌为人吃。又盛怒。曰。﹁臣口不能言。然臣期期知其不可。陛下虽欲废太子。臣期期不奉诏。﹂上欣然而笑。既罢。吕后侧耳於东箱听。见周昌。为跪谢曰。﹁微君。太子几废。﹂\\
\hspace*{2em}是後戚姬子如意为赵王。年十岁。高祖忧即万岁之後不全也。赵尧年少。为符玺御史。赵人方与公谓御史大夫周昌曰。﹁君之史赵尧。年虽少。然奇才也。君必异之。是且代君之位。﹂周昌笑曰。﹁尧年少。刀笔吏耳。何能至是乎。﹂居顷之。赵尧侍高祖。高祖独心不乐。悲歌。群臣不知上之所以然。赵尧进请问曰。﹁陛下所为不乐。非为赵王年少而戚夫人与吕后有卻邪。备万岁之後而赵王不能自全乎。﹂高祖曰。﹁然。吾私忧之。不知所出。﹂尧曰。﹁陛下独宜为赵王置贵彊相。及吕后。太子。群臣素所敬惮乃可。﹂高祖曰。﹁然。吾念之欲如是。而群臣谁可者。﹂尧曰。﹁御史大夫周昌。其人坚忍质直。且自吕后。太子及大臣皆素敬惮之。独昌可。﹂高祖曰。﹁善。﹂於是乃召周昌。谓曰。﹁吾欲固烦公。公彊为我相赵王。﹂周昌泣曰。﹁臣初起从陛下。陛下独柰何中道而弃之於诸侯乎。﹂高祖曰。﹁吾极知其左迁。然吾私忧赵王。念非公无可者。公不得已彊行。﹂於是徙御史大夫周昌为赵相。\\
\hspace*{2em}既行久之。高祖持御史大夫印弄之。曰。﹁谁可以为御史大夫者。﹂孰视赵尧。曰。﹁无以易尧。﹂遂拜赵尧为御史大夫。尧亦前有军功食邑。及以御史大夫从击陈豨有功。封为江邑侯。\\
\hspace*{2em}高祖崩。吕太后使使召赵王。其相周昌令王称疾不行。使者三反。周昌固为不遣赵王。於是高后患之。乃使使召周昌。周昌至。谒高后。高后怒而骂周昌曰。﹁尔不知我之怨戚氏乎。而不遣赵王。何。﹂昌既徵。高后使使召赵王。赵王果来。至长安月馀。饮药而死。周昌因谢病不朝见。三岁而死。\\
\hspace*{2em}後五岁。高后闻御史大夫江邑侯赵尧高祖时定赵王如意之画。乃抵尧罪。以广阿侯任敖为御史大夫。\\
\hspace*{2em}任敖者。故沛狱吏。高祖尝辟吏。吏系吕后。遇之不谨。任敖素善高祖。怒。击伤主吕后吏。及高祖初起。敖以客从为御史。守丰二岁。高祖立为汉王。东击项籍。敖迁为上党守。陈豨反时。敖坚守。封为广阿侯。食千八百户。高后时为御史大夫。三岁免。以平阳侯曹窋为御史大夫。高后崩。与大臣共诛吕禄等。免。以淮南相张苍为御史大夫。\\
\hspace*{2em}苍与绛侯等尊立代王为孝文皇帝。四年。丞相灌婴卒。张苍为丞相。\\
\hspace*{2em}自汉兴至孝文二十馀年。会天下初定。将相公卿皆军吏。张苍为计相时。绪正律历。以高祖十月始至霸上。因故秦时本以十月为岁首。弗革。推五德之运。以为汉当水德之时。尚黑如故。吹律调乐。入之音声。及以比定律令。若百工。天下作程品。至於为丞相。卒就之。故汉家言律历者。本之张苍。苍本好书。无所不观。无所不通。而尤善律历。\\
\hspace*{2em}张苍德王陵。王陵者。安国侯也。及苍贵。常父事王陵。陵死後。苍为丞相。洗沐。常先朝陵夫人上食。然后敢归家。\\
\hspace*{2em}苍为丞相十馀年。鲁人公孙臣上书言汉土德时。其符有黄龙当见。诏下其议张苍。张苍以为非是。罢之。其後黄龙见成纪。於是文帝召公孙臣以为博士。草土德之历制度。更元年。张丞相由此自绌。谢病称老。苍任人为中候。大为奸利。上以让苍。苍遂病免。苍为丞相十五岁而免。孝景前五年。苍卒。谥为文侯。子康侯代。八年卒。子类代为侯。八年。坐临诸侯丧後就位不敬。国除。\\
\hspace*{2em}初。张苍父长不满五尺。及生苍。苍长八尺馀。为侯。丞相。苍子复长。及孙类。长六尺馀。坐法失侯。苍之免相後。老。口中无齿。食乳。女子为乳母。妻妾以百数。尝孕者不复幸。苍年百有馀岁而卒。\\
\hspace*{2em}申屠丞相嘉者。梁人。以材官蹶张从高帝击项籍。迁为队率。从击黥布军。为都尉。孝惠时。为淮阳守。孝文帝元年。举故吏士二千石从高皇帝者。悉以为关内侯。食邑二十四人。而申屠嘉食邑五百户。张苍已为丞相。嘉迁为御史大夫。张苍免相。孝文帝欲用皇后弟窦广国为丞相。曰。﹁恐天下以吾私广国。﹂广国贤有行。故欲相之。念久之不可。而高帝时大臣又皆多死。馀见无可者。乃以御史大夫嘉为丞相。因故邑封为故安侯。\\
\hspace*{2em}嘉为人廉直。门不受私谒。是时太中大夫邓通方隆爱幸。赏赐累巨万。文帝尝燕饮通家。其宠如是。是时丞相入朝。而通居上傍。有怠慢之礼。丞相奏事毕。因言曰。﹁陛下爱幸臣。则富贵之。至於朝廷之礼。不可以不肃。﹂上曰。﹁君勿言。吾私之。﹂罢朝坐府中。嘉为檄召邓通诣丞相府。不来。且斩通。通恐。入言文帝。文帝曰。﹁汝第往。吾今使人召若。﹂通至丞相府。免冠。徒跣。顿首谢。嘉坐自如。故不为礼。责曰。﹁夫朝廷者。高皇帝之朝廷也。通小臣。戏殿上。大不敬。当斩。吏今行斩之。﹂通顿首。首尽出血。不解。文帝度丞相已困通。使使者持节召通。而谢丞相曰。﹁此吾弄臣。君释之。﹂邓通既至。为文帝泣曰。﹁丞相几杀臣。﹂\\
\hspace*{2em}嘉为丞相五岁。孝文帝崩。孝景帝即位。二年。晁错为内史。贵幸用事。诸法令多所请变更。议以谪罚侵削诸侯。而丞相嘉自绌所言不用。疾错。错为内史。门东出。不便。更穿一门南出。南出者。太上皇庙堧垣。嘉闻之。欲因此以法错擅穿宗庙垣为门。奏请诛错。错客有语错。错恐。夜入宫上谒。自归景帝。至朝。丞相奏请诛内史错。景帝曰。﹁错所穿非真庙垣。乃外堧垣。故他官居其中。且又我使为之。错无罪。﹂罢朝。嘉谓长史曰。﹁吾悔不先斩错。乃先请之。为错所卖。﹂至舍。因欧血而死。谥为节侯。子共侯蔑代。三年卒。子侯去病代。三十一年卒。子侯臾代。六岁。坐为九江太守受故官送有罪。国除。\\
\hspace*{2em}自申屠嘉死之後。景帝时开封侯陶青。桃侯刘舍为丞相。及今上时。柏至侯许昌。平棘侯薛泽。武彊侯庄青翟。高陵侯赵周等为丞相。皆以列侯继嗣。娖娖廉谨。为丞相备员而已。无所能发明功名有著於当世者。\\
\hspace*{2em}太史公曰。张苍文学律历。为汉名相。而绌贾生。公孙臣等言正朔服色事而不遵。明用秦之颛顼历。何哉。周昌。木彊人也。任敖以旧德用。申屠嘉可谓刚毅守节矣。然无术学。殆与萧。曹。陈平异矣。\\
\hspace*{2em}孝武时丞相多甚。不记。莫录其行起居状略。且纪征和以来。\\
\hspace*{2em}有车丞相。长陵人也。卒而有韦丞相代。韦丞相贤者。鲁人也。以读书术为吏。至大鸿胪。有相工相之。当至丞相。有男四人。使相工相之。至第二子。其名玄成。相工曰。﹁此子贵。当封。﹂韦丞相言曰。﹁我即为丞相。有长子。是安从得之。﹂後竟为丞相。病死。而长子有罪论。不得嗣。而立玄成。玄成时佯狂。不肯立。竟立之。有让国之名。後坐骑至庙。不敬。有诏夺爵一级。为关内侯。失列侯。得食其故国邑。韦丞相卒。有魏丞相代。\\
\hspace*{2em}魏丞相相者。济阴人也。以文吏至丞相。其人好武。皆令诸吏带剑。带剑前奏事。或有不带剑者。当入奏事。至乃借剑而敢入奏事。其时京兆尹赵君。丞相奏以免罪。使人执魏丞相。欲求脱罪而不听。复使人胁恐魏丞相。以夫人贼杀待婢事而私独奏请验之。发吏卒至丞相舍。捕奴婢笞击问之。实不以兵刃杀也。而丞相司直繁君奏京兆尹赵君迫胁丞相。诬以夫人贼杀婢。发吏卒围捕丞相舍。不道。又得擅屏骑士事。赵京兆坐要斩。又有使掾陈平等劾中尚书。疑以独擅劫事而坐之。大不敬。长史以下皆坐死。或下蚕室。而魏丞相竟以丞相病死。子嗣。後坐骑至庙。不敬。有诏夺爵一级。为关内侯。失列侯。得食其故国邑。魏丞相卒。以御史大夫邴吉代。\\
\hspace*{2em}邴丞相吉者。鲁国人也。以读书好法令至御史大夫。孝宣帝时。以有旧故。封为列侯。而因为丞相。明於事。有大智。後世称之。以丞相病死。子显嗣。後坐骑至庙。不敬。有诏夺爵一级。失列侯。得食故国邑。显为吏至太仆。坐官秏毛乱。身及子男有奸赃。免为庶人。\\
\hspace*{2em}邴丞相卒。黄丞相代。长安中有善相工田文者。与韦丞相。魏丞相。邴丞相微贱时会於客家。田文言曰。﹁今此三君者。皆丞相也。﹂其後三人竟更相代为丞相。何见之明也。\\
\hspace*{2em}黄丞相霸者。淮阳人也。以读书为吏。至颍川太守。治颍川。以礼义条教喻告化之。犯法者。风晓令自杀。化大行。名声闻。孝宣帝下制曰。﹁颍川太守霸。以宣布诏令治民。道不拾遗。男女异路。狱中无重囚。赐爵关内侯。黄金百斤。﹂徵为京兆尹而至丞相。复以礼义为治。以丞相病死。子嗣。後为列侯。黄丞相卒。以御史大夫于定国代。于丞相已有廷尉传。在张廷尉语中。于丞相去。御史大夫韦玄成代。\\
\hspace*{2em}韦丞相玄成者。即前韦丞相子也。代父。後失列侯。其人少时好读书。明於诗。论语。为吏至卫尉。徙为太子太傅。御史大夫薛君免。为御史大夫。于丞相乞骸骨免。而为丞相。因封故邑为扶阳侯。数年。病死。孝元帝亲临丧。赐赏甚厚。子嗣後。其治容容随世俗浮沈。而见谓谄巧。而相工本谓之当为侯代父。而後失之。复自游宦而起。至丞相。父子俱为丞相。世间美之。岂不命哉。相工其先知之。韦丞相卒。御史大夫匡衡代。\\
\hspace*{2em}丞相匡衡者。东海人也。好读书。从博士受诗。家贫。衡佣作以给食饮。才下。数射策不中。至九。乃中丙科。其经以不中科故明习。补平原文学卒史。数年。郡不尊敬。御史徵之。以补百石属荐为郎。而补博士。拜为太子少傅。而事孝元帝。孝元好诗。而迁为光禄勋。居殿中为师。授教左右。而县官坐其旁听。甚善之。日以尊贵。御史大夫郑弘坐事免。而匡君为御史大夫。岁馀。韦丞相死。匡君代为丞相。封乐安侯。以十年之间。不出长安城门而至丞相。岂非遇时而命也哉。\\
\hspace*{2em}太史公曰。深惟士之游宦所以至封侯者。微甚。然多至御史大夫即去者。诸为大夫而丞相次也。其心冀幸丞相物故也。或乃阴私相毁害。欲代之。然守之日久不得。或为之日少而得之。至於封侯。真命也夫。御史大夫郑君守之数年不得。匡君居之未满岁。而韦丞相死。即代之矣。岂可以智巧得哉。多有贤圣之才。困戹不得者众甚也。\\

% ===== 索隐述赞（司马贞《史记索隐》）· 统一样式 =====
% 题头「索隐述赞曰」五字并入反白框；正文每四句一列，定宽盒 \hfil 居列中；
% \fontsize{24}{41} 列距对齐全书界行（121+41k）；\clearpage 使述赞独占一页。
\clearpage
\shuzanhead{索隐述赞曰}\linebreak
{\shuzanbody\fontsize{24}{41}\selectfont
\hbox to \hsize{\hfil 张苍主计。天下作程。孙臣始绌。秦历尚行。\hfil}\linebreak
\hbox to \hsize{\hfil 御史亚相。相国阿衡。申屠面折。周子廷争。\hfil}\linebreak
\hbox to \hsize{\hfil 其他娖々。无所发明。\hfil}%
\par}
\end{document}
